\subsection{Variational inference and the ELBO}
\label{subsec:vi-elbo}

\subsubsection{From EM free energy to the ELBO}

The transition from EM is literal rather than metaphorical.
\Cref{prop:em-free-energy-coordinate-ascent} introduced the free-energy functional for a latent-variable model \(p_\theta(x,z)\):
\[
  \mathcal{F}(q,\theta;x)
  =
  \E_q[\log p_\theta(x,Z)] + H(q),
\]
and showed that exact EM alternates between the unrestricted update \(q(z)=p_\theta(z\mid x)\) and a maximization in \(\theta\).
Variational inference keeps the same objective but restricts the latent law \(q\) to a tractable family.
This is the only conceptual change: the model \(p_\theta(x,z)\) stays in place, while the exact E-step is relaxed into an optimization problem.

For this subsection we regard \(\theta\) as fixed and write
\[
  p_\theta(x,z) = p_\theta(x\mid z)\,p_\theta(z),
\]
so the posterior is
\[
  p_\theta(z\mid x) = \frac{p_\theta(x,z)}{p_\theta(x)},
  \qquad
  p_\theta(x) = \int p_\theta(x,z)\,dz.
\]
The evidence \(p_\theta(x)\) is exactly the quantity that makes direct Bayesian inference difficult:
it is the normalizing constant of the posterior and typically requires a high-dimensional integral.
Modern latent-variable methods usually introduce the ELBO at exactly this point; see \citet[\S2.2]{kingma_welling_2014_aevb} for the VAE version of the same identity.

With \(\theta\) fixed, the free-energy functional becomes
\[
  \mathcal{F}(q,\theta;x)
  \coloneqq
  \E_q\!\bigl[\log p_\theta(x,z)\bigr]
  +
  H(q),
\]
where $H(q)=-\E_q[\log q(z)]$ is the entropy.
Variational inference keeps this same objective but restricts the $q$-update from all densities to a smaller family
\[
  \mathcal{Q} = \{q_\lambda(z): \lambda\in\Lambda\}
\]
The ELBO is simply the computational form of this restricted free-energy problem.

\subsubsection{ELBO--KL identity}

This is the same decomposition that underlies EM, now viewed with \(\theta\) fixed and \(q\) ranging over a
restricted family. Writing it explicitly fixes the KL direction and makes clear why the ELBO is the
computational surrogate for posterior approximation.

\begin{proposition}[Free-energy and ELBO identity]
\label{prop:elbo-identity}
For any density \(q(z)\) supported where \(p_\theta(x,z)>0\),
\begin{equation}
  \log p_\theta(x)
  =
  \mathcal{F}(q,\theta;x)
  +
  \operatorname{KL}\!\bigl(q(z)\,\|\,p_\theta(z\mid x)\bigr),
  \label{eq:elbo-identity}
\end{equation}
Consequently,
\[
  \mathcal{F}(q,\theta;x) \le \log p_\theta(x),
\]
and maximizing \(\mathcal{F}(q,\theta;x)\) over the variational family is equivalent to minimizing \(\operatorname{KL}(q\|p_\theta(z\mid x))\).
\end{proposition}
\begin{proof}
Starting from the definition of KL divergence and substituting the posterior density,
\[
  \operatorname{KL}\!\bigl(q(z)\,\|\,p_\theta(z\mid x)\bigr)
  =
  \int q(z)\log\frac{q(z)}{p_\theta(x,z)/p_\theta(x)}\,dz.
\]
Expanding the logarithm gives
\[
  \operatorname{KL}\!\bigl(q(z)\,\|\,p_\theta(z\mid x)\bigr)
  =
  \int q(z)\log q(z)\,dz
  -
  \int q(z)\log p_\theta(x,z)\,dz
  +
  \log p_\theta(x).
\]
Rearranging yields
\[
  \log p_\theta(x)
  =
  \E_q\!\bigl[\log p_\theta(x,z)-\log q(z)\bigr]
  +
  \operatorname{KL}\!\bigl(q(z)\,\|\,p_\theta(z\mid x)\bigr),
\]
which is exactly \eqref{eq:elbo-identity} because
\[
  \mathcal{F}(q,\theta;x)
  =
  \E_q[\log p_\theta(x,z)] + H(q)
  =
  \E_q[\log p_\theta(x,z)-\log q(z)].
\]
Since KL divergence is nonnegative, \eqref{eq:elbo-identity} implies \(\mathcal{F}(q,\theta;x)\le \log p_\theta(x)\), and the difference between the two objectives is the constant \(\log p_\theta(x)\).%
\end{proof}

It is conventional to rename the same quantity
\begin{equation}
  \mathcal{L}(q;\theta,x)
  \coloneqq
  \E_q\!\bigl[\log p_\theta(x,z)-\log q(z)\bigr]
  =
  \mathcal{F}(q,\theta;x)
  \label{eq:elbo-definition}
\end{equation}
and call it the \emph{evidence lower bound}, or ELBO.
The new notation emphasizes the computational point: after fixing \(\theta\), the objective depends only on expectations under the tractable law \(q\), while the intractable evidence \(\log p_\theta(x)\) appears only as an additive constant.
Thus
\begin{equation}
  q^\star
  \in
  \argmax_{q\in\mathcal{Q}} \mathcal{L}(q;\theta,x)
  \quad\Longleftrightarrow\quad
  q^\star
  \in
  \argmin_{q\in\mathcal{Q}} \operatorname{KL}\!\bigl(q(z)\,\|\,p_\theta(z\mid x)\bigr).
  \label{eq:vi-kl-objective}
\end{equation}
This is where VI diverges from both EM and MCMC.
EM keeps the exact posterior update but alternates it with a parameter maximization, while MCMC constructs a chain with the posterior as its stationary law.
VI instead replaces posterior sampling or exact posterior evaluation by deterministic optimization over a restricted family.

\subsubsection{Reconstruction--regularization decomposition}

The ELBO also admits a decomposition that is often more interpretable.
Using \(p_\theta(x,z)=p_\theta(x\mid z)p_\theta(z)\) in \eqref{eq:elbo-definition} gives
\begin{equation}
  \mathcal{L}(q;\theta,x)
  =
  \E_q[\log p_\theta(x\mid z)]
  -
  \operatorname{KL}\!\bigl(q(z)\,\|\,p_\theta(z)\bigr).
  \label{eq:elbo-reconstruction-regularization}
\end{equation}
The first term rewards latent configurations that explain the data well, while the second penalizes variational distributions that drift too far from the prior.
In latent-variable models this is the usual ``fit minus regularization'' interpretation.
The decoder/prior language from variational autoencoders is exactly this decomposition written for a model with a fixed prior and a learned conditional distribution \(p_\theta(x\mid z)\) \citep[\S2.2]{kingma_welling_2014_aevb}.

\subsubsection{What forward KL tends to prefer}

There is, however, an important geometric asymmetry in \eqref{eq:vi-kl-objective}.
The direction $\operatorname{KL}(q\|p)$ heavily penalizes placing mass where the posterior is small, but it is more tolerant of failing to cover every region where the posterior is large.
With a restrictive family this often leads, as a heuristic rule of thumb, to approximations that understate posterior variance or choose one mode of a multimodal posterior rather than averaging across all of them.
See \citet[\S2.4 and Fig.~1]{blei_kucukelbir_mcauliffe_2017_vi_review} for a concrete mean-field example of this variance-underestimation effect.
This is not a flaw in the ELBO identity itself.
It is the geometric tradeoff created by combining an asymmetric divergence with a tractable but limited family.
The next subsection makes that tradeoff concrete by studying the most classical restricted family, namely a fully factorized mean-field approximation.
